\section{Conclusion}
\label{sec:ch7-conclusion}
This study examined whether the proceduralization of the cooperation and consistency mechanism introduced by the Procedural Regulation would curb the discretion of the \gls{lsa} in composite cross-border cases.
To answer this question, it analyzed the functioning of the \gls{gdpr}'s integrated system of cooperation and identified new opportunities for peer pressure in the horizontal and vertical interactions between \glspl{dpa}, situating these against the theoretical framework of the duties of care and to give reasons.
Following a doctrinal analysis of the principle of good administration, and of these two duties in particular, it mapped the deficiencies of the previous system of cooperation and asked whether, in those cases where violations of the two duties had occurred, the Procedural Regulation's structure of procedural steps and deadlines would have produced a different outcome.

The study concluded that the Procedural Regulation does introduce formal tools that, if employed, may help expand the scope of investigations and provide a check on the \gls{lsa}'s discretion.
By specifying multiple aspects of the cooperation procedure, it now offers \glspl{sa} a structured framework to follow in cross-border enforcement.
Yet while the uncertainty of the pre-reform system may have contributed to the procedural gaps underlying these violations, the reformed flow of cooperation does not address the more profound problems in the \gls{gdpr}'s enforcement system, such as the chronic underfunding of \glspl{dpa} and the divergence of national enforcement priorities and procedures.
Informational asymmetries may therefore persist where \glspl{sa} lack the resources needed to make use of these formal tools.
In some cases, moreover, the Regulation provides procedures that are likely to reduce enforcement, most notably in the case of early resolution.
And the intricate scheme of procedural deadlines may prove, in practice, to be a mere organizational benchmark rather than an enforceable procedural requirement.

Failure to protect the fundamental right to data protection through public enforcement may endanger not only that right but the consistent application of the \gls{gdpr} as a whole, and with it the protection of other fundamental rights such as freedom of expression, non-discrimination, and the right to private life.

\subsection{Contribution}
\label{sec:ch7-contribution}
This thesis has offered the first structured assessment of the Procedural Regulation against the duties of care and to give reasons, understood as sub-principles of the principle of good administration.
It has done so by conceptualizing the \gls{gdpr}'s integrated enforcement system as a form of cooperation between peers who may influence one another, through the use of procedural tools, towards more carefully investigated and better reasoned enforcement outcomes.

\subsection{Areas for Future Research}
\label{sec:ch7-future-research}
This study has concentrated on the horizontal interactions between \glspl{dpa} and on their vertical relationship with the \gls{edpb}.
The Procedural Regulation, however, also introduces new duties for \glspl{dpa} towards the parties under investigation and the complainants, a dimension that future work could explore.
The analysis has likewise been limited to two sub-elements of the principle of good administration, and more specifically to their objective dimension.
A more extensive study could address both the objective and subjective dimensions of the principle, as well as other principles of \gls{eu} law, such as subsidiarity, proportionality, sincere cooperation, and the right to an effective remedy.
Finally, the Regulation has not yet entered into force, and much of its practical impact can only truly be assessed once the first cross-border cases have been processed under its provisions.